% ECONOMICS_PROBLEM: {"id": "F38-60", "title": "Optimal capital-flow policy mix", "jel_code": "F38", "source_id": "OP-0337"}
% !TeX program = xelatex
\documentclass[11pt,a4paper]{article}
\usepackage[margin=23mm]{geometry}
\usepackage{fontspec}
\setmainfont{Latin Modern Roman}
\usepackage{amsmath,amssymb,mathtools}
\usepackage{xcolor,enumitem,booktabs,longtable,array,xurl,hyperref}
\definecolor{muted}{HTML}{53636C}
\hypersetup{colorlinks=true,urlcolor=blue,linkcolor=blue}
\setlength{\parindent}{0pt}
\setlength{\parskip}{5pt}
\newcommand{\R}{\mathbb R}
\newcommand{\E}{\mathbb E}
\newcommand{\N}{\mathbb N}
\newcommand{\ind}{\mathbf 1}
\newcommand{\Var}{\operatorname{Var}}
\newcommand{\Cov}{\operatorname{Cov}}
\newcommand{\supp}{\operatorname{supp}}
\DeclareMathOperator*{\argmin}{arg\,min}
\DeclareMathOperator*{\argmax}{arg\,max}
\newcommand{\entrymeta}[3]{{\sffamily\small\color{muted}\textbf{\texttt{#1}}\quad|\quad #2\par #3\par}}
\newcommand{\Problemref}[1]{\texttt{#1}}
\newcommand{\JELref}[2]{\texttt{#2} (JEL \texttt{#1})}
\begin{document}
\section*{Optimal capital-flow policy mix}
\textbf{Problem ID:} \texttt{F38-60}\quad \textbf{JEL:} \texttt{F38}\par
% BEGIN ECONOMICS STATEMENT
\label{problem:OP-0337}
{\sffamily\small\color{muted}Original subject group: International finance and exchange rates\par}
\label{definitions:problem:30007619}
\textbf{Audit classification:} Published research agenda. The Bank agenda supports combining international tools; the IMF source distinguishes fundamental from financial currency shocks. This exact four-instrument robust rule is editorial, with commitment and discretion separated.
\subsection*{Definitions and precise research target}
Fix an open economy with sticky prices, foreign-currency liabilities and financial-market frictions specified as primitives. At date \(t\), authorities observe state \(x_t\) and choose an instrument vector \(a_t=(r_t,f_t,c_t,m_t)\): interest-rate policy, foreign-exchange intervention, capital-flow measures and macroprudential restrictions. Here these symbols denote policy controls, not asset returns or consumption. Feasible rules \(a_t=\pi(x_t)\) obey reserve availability, legal limits and information constraints. Let \(W(\pi;P)\) be expected discounted household welfare under equilibrium private responses and shock law \(P\). The task is to identify or characterize
\[J(\pi)=\inf_{P\in\mathcal P}W(\pi;P),\qquad W_*=\sup_{\pi\in\Pi}J(\pi),\]
where \(\Pi\) is a stated feasible class and \(\mathcal P\) captures uncertainty about frictions and shocks. Derive conditions under which exchange-rate adjustment alone is optimal and when additional instruments improve welfare. Separate shocks to fundamentals from impaired currency-market functioning; the same observed depreciation can imply different optimal responses. Include intervention sterilization, balance-sheet effects and policy interactions. Identify the sufficient statistics needed to select tools in real time and quantify losses from misclassification. A policy combination optimal in one calibrated economy cannot be asserted optimal for every institutional setting.

\textbf{Additional formal definitions and scope (editorial).} The state \(x_t\) includes reserve assets, currency liabilities, inflation and any information used by the authority. Interest-rate control \(r_t\), FX purchases \(f_t\), capital-flow wedge \(c_t\) and macroprudential wedge \(m_t\) each have declared units, timing and implementation equations. Specify sterilization and the government/central-bank budget so intervention feasibility is operational. Feasible \(\pi\) are nonanticipating measurable rules; if private expectations depend on future interventions, state whether the authority commits or reoptimizes. The uncertainty set includes shock laws and structural friction parameters, and privately anticipated policies must be re-solved for each. Finding an optimum in a known finite-state model is a standard control problem; the wider question concerns justified frictions, identification of tool effects and robustness. Commitment and discretionary optimization are different variants. For a nonempty feasible set and a finite optimal value, the design target is an \(\varepsilon\)-optimal feasible rule for each specified \(\varepsilon>0\): its cost is at most the infimum plus \(\varepsilon\), or its welfare is at least the supremum minus \(\varepsilon\). Report infeasibility or an unbounded value separately. An exact optimizer is requested only under additional attainment assumptions; it need not exist for the general rule class.
\subsection*{Variants and assumption changes}
\begin{enumerate}
\item \textbf{One identified friction, commitment (editorial).} Specify a finite-state model and optimize feasible FX/interest-rate rules with a fixed capital-flow regime.
\item \textbf{Joint instruments under uncertainty (editorial).} Allow currency mismatch, sticky prices and multiple tools over \(\mathcal P\); compare commitment and discretionary rules using the same welfare accounting.
\end{enumerate}
\subsection*{References and source audit}
\begin{enumerate}
\item Official January 9, 2026 web version; locator: Interconnected World / International finance, combination-of-tools paragraph. Broad agenda; exact model editorial. URL: \url{https://www.bankofengland.co.uk/research/bank-of-england-agenda-for-research}.
\item Official IMF/elibrary indexed records verify five authors, 2026 SDN/2026/003 and DOI. Official release is September 17, 2026; URL contains September 16. Locator: Official indexed abstract and currency-premium/decomposition passages; no exact full-page locator certified. Support: Fundamental/financial FX identification and policy background. Landing page initially opened; later full-text fetches failed. Indexed primary passages inspected only. URL: \url{https://www.imf.org/en/publications/staff-discussion-notes/issues/2026/09/16/drivers-of-exchange-rates-in-emdes-implications-for-foreign-exchange-intervention-578715}.
\end{enumerate}
\begin{enumerate}\raggedright
\item\hypertarget{definitions:source:30007619:1}{} Bank of England. 2026. \emph{Bank of England Agenda for Research, 2025–2028}. Interconnected World / International finance, combination-of-tools paragraph.
\url{https://www.bankofengland.co.uk/research/bank-of-england-agenda-for-research}.
\item\hypertarget{definitions:source:30007619:2}{} Ece Ozge Emeksiz; Andres Fernandez; Nikhil Patel; Ivan Petrella; Tatjana Schulze. 2026. \emph{Drivers of Exchange Rates in EMDEs: Implications for Foreign Exchange Intervention}. Official indexed abstract and currency-premium/decomposition passages; no exact full-page locator certified.
\url{https://www.imf.org/en/publications/staff-discussion-notes/issues/2026/09/16/drivers-of-exchange-rates-in-emdes-implications-for-foreign-exchange-intervention-578715}.
DOI: \url{https://doi.org/10.5089/9798229055208.006}.
\end{enumerate}

\subsection*{Common conventions: Source mathematical conventions}

These conventions apply to each entry unless explicitly replaced.

\textbf{Models and identification.} A primitive model \(m\) specifies all probability laws, feasible choices, information, equilibrium and measurement rules used by its target. The admissible class \(\mathcal M\) is fixed independently of outcomes. If \(P_O(m)\) is its observable probability law and \(\tau(m)\) the target, its identified set at observable law \(P\) is
\[
 \mathcal I_\tau(P)=\{\tau(m):m\in\mathcal M,\ P_O(m)=P\}.
\]
Point identification means that this set is a singleton. An empty set signals incompatibility of data and assumptions. Coordinate bounds need not describe the joint set. Bounds are sharp only if every retained target is attainable or arbitrarily approachable by compatible models. Finite bounds require explicit restrictions on unobserved quantities; unrestricted confounding, hidden wealth, extrapolation or arbitrary equilibrium selection can yield uninformative sets.

\textbf{Causal interventions.} A potential outcome \(Y_i(p)\) is the outcome under a complete policy regime \(p\), including timing, financing, information and market spillovers. The target population and units are fixed before treatment. With interference, use the complete assignment vector or a justified exposure mapping; individual no-interference notation alone cannot identify an economy-wide intervention. A difference of mean potential outcomes does not require the cross-world joint distribution. Conditioning on post-treatment migration, survival, enrollment or employment changes the estimand and needs separate justification.

\textbf{Statistical statements.} An estimator, confidence set, forecast or test specifies a sampling law, model class and loss/coverage criterion. Uniform \((1-\alpha)\)-coverage means \(\inf_{m\in\mathcal M}\Pr_m\{\tau(m)\in C_n\}\geq1-\alpha\), or an explicitly asymptotic version. The whole parameter space trivially has coverage, so informativeness requires a stated width, risk or power objective. A forecasting improvement is relative to a named benchmark, information set and evaluation population.

\textbf{Equilibrium and optimization.} An equilibrium comprises feasible plans, optimality, beliefs where needed, budget constraints and all market/resource clearing. The primitives or a stated benchmark define each correspondence \(\mathcal E(p;m)\). Nonemptiness is assumed or proved, never inferred just from compact policy sets. A selected-equilibrium target uses a declared selection rule; a robust target quantifies over all permitted equilibria. Use suprema or infima unless attainment is established. An \(\varepsilon\)-optimal policy is feasible and within \(\varepsilon\) of the finite optimum value. Report an empty feasible set separately.

\textbf{Welfare and accounting.} Normative utility, interpersonal normalization, social weights, numeraire, discounting and the use of public receipts are primitives. Behavioral utility need not coincide with the welfare criterion. Household welfare includes ownership income and rents once; adding profits again to compensating variation generally double-counts them. A monetary equivalent is defined by an explicit transfer and price convention, with monotonicity and range conditions. Average effects, mechanism contributions, transfers and welfare losses are different targets. Infinite sums require convergence and dynamic feasible plans require terminal or no-Ponzi conditions.

\textbf{Source versions.} A working-paper date, revision date and journal publication year are distinguished. A DOI for a journal version is not automatically the DOI of an earlier manuscript. Locators refer to the stated version and printed pages where specified. A metadata-only check does not verify a claim about a full-text conclusion. Every entry's source subsection narrows the provenance claim accordingly.
% END ECONOMICS STATEMENT
\end{document}
